\section{Radio Access Configuration}
\label{sec:radio}

The simulated deployment targets 5G NR Band~n78 at $f_c = 3.6\,\text{GHz}$
(primary Sub-6\,GHz micro-cell allocation in 3GPP~TS~38.104~\cite{3gpp_ts38104,3gpp_nr_channel}).
All nine rooftop sites share the RF and antenna configuration in Table~\ref{tab:radio_spec}.

\begin{table}[htbp]
  \caption{Radio link and OFDM configuration (per BS site)}
  \label{tab:radio_spec}
  \centering
  \begin{tabular}{ll}
    \toprule
    Parameter & Value \\
    \midrule
    Operating band          & 5G NR n78 ($3.3$--$3.8\,\text{GHz}$) \\
    Carrier frequency       & $f_c = 3.6\,\text{GHz}$ \\
    Subcarrier spacing      & $\Delta f = 30\,\text{kHz}$ (numerology $\mu=1$) \\
    FFT size / subcarriers  & $N_{\mathrm{FFT}} = 3168$ \\
    Transmit power          & $+40\,\text{dBm}$ per site \\
    BS antenna array        & $4\times4$ planar, 16 elements,
                              TR~38.901 pattern \\
    BS antenna height       & $19\,\text{m}$ \\
    UE antenna              & Isotropic, 1 element \\
    UE antenna height       & $1.25\,\text{m}$ \\
    Max reflection depth    & 4 \\
    Max ray depth           & 4 \\
    \bottomrule
  \end{tabular}
\end{table}

The BS uses a $4\times4$ dual-polarised TR~38.901 planar array (16 elements, half-wavelength spacing)
for angular diversity and cross-BS UE fingerprint distinguishability~\cite{hoydis2023sionna_rt,dahlman2018_5gnr}.
The UE is modelled with a single isotropic antenna.
Ray tracing retains the 30 strongest paths per BS--UE link (up to 4 reflection interactions).
The OFDM transfer function is evaluated on the full $N_{\mathrm{FFT}}=3168$ subcarrier grid
and downsampled for CSI fingerprints in Sections~\ref{sec:localization}--\ref{sec:channel_charting}.
